\documentclass[12pt,a4paper]{article}

% ─── Codificação e língua ────────────────────────────────────────────────────
\usepackage[utf8]{inputenc}
\usepackage[T1]{fontenc}
\usepackage[brazil]{babel}

% ─── Matemática ─────────────────────────────────────────────────────────────
\usepackage{amsmath, amssymb, amsthm, mathrsfs}

% ─── Layout ─────────────────────────────────────────────────────────────────
\usepackage[top=2.5cm, bottom=2.5cm, left=2.8cm, right=2.8cm]{geometry}
\usepackage{setspace}
\onehalfspacing

% ─── Tabelas e floats ────────────────────────────────────────────────────────
\usepackage{booktabs}
\usepackage{array}
\usepackage{float}
\usepackage{enumerate}

% ─── Hiperlinks ─────────────────────────────────────────────────────────────
\usepackage[colorlinks=true, linkcolor=black, citecolor=black, urlcolor=blue]{hyperref}

% ─── Código ─────────────────────────────────────────────────────────────────
\usepackage{listings}
\lstset{basicstyle=\ttfamily\small, frame=single, breaklines=true}

% ─── Ambientes de teorema  (mesmo padrão dos papers [1] e [2]) ───────────────
\newtheorem{theorem}{Teorema}[section]
\newtheorem{lemma}[theorem]{Lema}
\newtheorem{proposition}[theorem]{Proposição}
\newtheorem{corollary}[theorem]{Corolário}
\newtheorem{conjecture}[theorem]{Conjectura}
\newtheorem{hypothesis}[theorem]{Hipótese}

\theoremstyle{definition}
\newtheorem{definition}[theorem]{Definição}
\newtheorem{example}[theorem]{Exemplo}

\theoremstyle{remark}
\newtheorem{remark}[theorem]{Observação}

% ─── Macros ──────────────────────────────────────────────────────────────────
\newcommand{\NN}{\mathbb{N}}
\newcommand{\ZZ}{\mathbb{Z}}
\newcommand{\RR}{\mathbb{R}}
\newcommand{\PP}{\mathbb{P}}          % conjunto dos primos
\newcommand{\A}{\mathcal{A}}         % conjunto dos pares
\newcommand{\Grid}{G_{3,N}}
\newcommand{\Om}{\Omega_N}           % peso oracle-based
\newcommand{\Oms}{\Omega^*_N}        % peso oracle-free
\newcommand{\Ss}{\mathfrak{S}}       % série singular
\newcommand{\Ku}{K_u}
\newcommand{\lnN}{\ln N}

% ─── Metadados ───────────────────────────────────────────────────────────────
\title{%
	O Invariante Âncora da Grade $G_{3,N}$ e um Estimador\\
	Geométrico Oracle-Free para a Conjectura de Goldbach%
}
\author{%
	Tiago Bandeira\\[4pt]
	\small\texttt{tiagobandeira.goldbach2025@gmail.com}%
}
\date{Maio de 2026\\[4pt]
	\small Preprint --- Paper~3 da série; complementa \cite{bandeira2026unif,bandeira2026soma}%
}

% ════════════════════════════════════════════════════════════════════════════
\begin{document}
	\maketitle
	
	% ─── Nota do autor ───────────────────────────────────────────────────────────
	\begin{center}
		\small\itshape
		\parbox{0.88\textwidth}{%
			\textbf{Nota do Autor.}
			Este trabalho constitui uma investigação exploratória e não reivindica
			uma demonstração da Conjectura de Goldbach.
			Resultados rigorosos são explicitamente identificados;
			hipóteses heurísticas são declaradas como tais.
			Comentários e críticas matemáticas são bem-vindos.%
		}
	\end{center}
	\bigskip
	
	% ─── Abstract ────────────────────────────────────────────────────────────────
	\begin{abstract}
		O peso geométrico $\Om(n)$, introduzido em \cite{bandeira2026unif} para medir
		o ``campo gravitacional primo'' ao redor de cada célula~$n$ da grade $\Grid$,
		é computável apenas com um oráculo de primalidade --- limitação que impede
		sua aplicação como crivo \emph{a priori}.
		Este trabalho resolve essa limitação explorando o \textbf{Invariante Âncora}
		demonstrado em \cite{bandeira2026soma}: em toda janela~$W_i$ a soma das três
		configurações canônicas é $S(i) = 3(2N+2i+1)$, e o elemento central da
		linha mediana satisfaz $f(2,i{+}1) = S(i)/3$.
		Definimos o complemento oracle-free $C(n,i) = S(i)-n$ e o peso
		\[
		\Oms(n) \;=\; \frac{1}{|W(n)|}\sum_{i\in W(n)}\Ss\!\bigl(C(n,i)\bigr),
		\qquad
		\Ss(C) \;=\!\!\prod_{\substack{p\mid C\\p>2}}\!\frac{p-1}{p-2},
		\]
		onde $\Ss(C)$ é a Série Singular de Hardy--Littlewood, computável da
		fatoração de~$C$ sem nenhum teste de primalidade de vizinhos.
		Verificação computacional em $\Grid$ para $N$ de $10^3$ a $10^5$ mostra que
		$\Oms$ discrimina estatisticamente $P_1$ (primos) de $P_2$ (semiprimos)
		\textbf{em todos os $N$ testados}, incluindo os $N$ com fatoração restrita
		a $\{2,3,5\}$ onde o peso original $\Om$ falha completamente.
		A correlação entre $\Om$ e $\Oms$ é próxima de zero ($r\approx-0.06$ para
		$P_1$), confirmando que os dois pesos capturam sinais aritmeticamente
		independentes.
		Discutimos as implicações para a Conjectura de Goldbach e para outros
		problemas clássicos sobre primos --- primos gêmeos, Legendre, Oppermann,
		Brocard --- e delineamos as direções de pesquisa abertas por este resultado.
	\end{abstract}
	
	\tableofcontents
	\bigskip
	
	% ════════════════════════════════════════════════════════════════════════════
	\section{Introdução}
	
	A grade $\Grid$, introduzida em \cite{bandeira2026unif}, organiza os
	primeiros $3N$ números ímpares em três linhas de $N$ colunas.
	Janelas $3\times3$ consecutivas --- denominadas $W_i$ --- contêm três
	configurações canônicas (diagonal principal~$D_1$, diagonal
	secundária~$D_2$ e coluna central~$C_v$).
	Em \cite{bandeira2026unif} demonstrou-se, via Borel--Cantelli sob
	quasi-independência de Cramér, que trios primos nessas configurações
	ocorrem quase certamente infinitas vezes, e cada trio primo testemunha
	simultaneamente três representações de Goldbach binária e uma de Goldbach
	ternária.
	
	O peso geométrico $\Om(n)$, também definido em \cite{bandeira2026unif},
	quantifica a fração média de vizinhos primos de~$n$ nas janelas que o
	contêm.
	Testes de Kolmogorov--Smirnov (KS) e Mann--Whitney (MW) mostraram que
	$\Om$ discrimina $P_1$ de $P_2$ para $N$ suficientemente grande ou livre
	de obstruções modulares.
	Entretanto, o cálculo de $\Om(n)$ requer saber se cada vizinho é primo ---
	um oráculo de primalidade.
	Isso impede sua aplicação como crivo \emph{a priori} e, como observado na
	Seção~11.4 de \cite{bandeira2026unif}, é precisamente o nó computacional
	do programa do crivo geométrico.
	
	Em \cite{bandeira2026soma} foi demonstrado o \textbf{Lema da Soma Constante}:
	em toda janela~$W_i$, as três configurações $D_1$, $D_2$ e $C_v$ têm sempre
	a mesma soma $S(i) = 3(2N+2i+1)$.
	Este resultado é puramente algébrico --- independe de primalidade.
	Uma consequência imediata, não explorada em \cite{bandeira2026soma}, é que
	o elemento central $f(2,i{+}1) = S(i)/3$ é um \textbf{invariante âncora}:
	ele determina, por posição, o complemento aritmético de qualquer outro
	elemento da janela.
	
	Este trabalho explora essa consequência.
	Definimos $C(n,i) = S(i)-n$ e construímos $\Oms(n)$ como a média da Série
	Singular $\Ss(C(n,i))$ sobre as janelas de~$n$.
	Esse peso é oracle-free: depende apenas da posição de~$n$ e da fatoração
	de~$C$, sem acessar a primalidade de qualquer vizinho.
	A verificação computacional revela que $\Oms$ supera $\Om$ em robustez,
	discriminando $P_1$ de $P_2$ uniformemente em todos os $N$ testados.
	
	% ════════════════════════════════════════════════════════════════════════════
	\section{Preliminares e Notação}
	
	Adotamos a notação unificada de \cite{bandeira2026unif}.
	
	\begin{definition}[Grade $G_{3,N}$]
		Fixado $N\in\NN$ com $N\geq3$, a grade $\Grid$ é uma matriz de três linhas e
		$N$ colunas cujas células $(r,c)$, $r\in\{1,2,3\}$, $c\in\{1,\dots,N\}$,
		recebem os ímpares pelo mapeamento linha a linha:
		\[
		f(r,c) = 2\bigl[(r-1)N+c\bigr]-1.
		\]
		A janela $W_i$, $i=1,\dots,N-2$, é a submatriz formada pelas colunas
		$i,\,i{+}1,\,i{+}2$.
		As configurações canônicas são: $D_1 = \{(1,i),(2,i{+}1),(3,i{+}2)\}$,
		$D_2 = \{(3,i),(2,i{+}1),(1,i{+}2)\}$, $C_v = \{(1,i{+}1),(2,i{+}1),(3,i{+}1)\}$.
	\end{definition}
	
	\begin{lemma}[Soma constante {\cite[Lema~1]{bandeira2026soma}}]
		\label{lem:soma}
		Para qualquer janela $W_i$ de $\Grid$:
		\[
		D_1 = D_2 = C_v = S(i) := 3(2N+2i+1).
		\]
	\end{lemma}
	
	\begin{corollary}[Elemento âncora]
		\label{cor:ancora}
		O elemento $f(2,i{+}1)$ da linha mediana satisfaz
		\[
		f(2,i{+}1) \;=\; 2N+2i+1 \;=\; \frac{S(i)}{3}.
		\]
	\end{corollary}
	
	\begin{proof}
		Direto de $f(2,i{+}1)=2[(2-1)N+(i{+}1)]-1=2N+2i+1$ e $S(i)/3=2N+2i+1$.
	\end{proof}
	
	\begin{remark}[Paridade do complemento]
		\label{rem:paridade}
		Para qualquer célula $n=f(r,c)$ e janela $i$ contendo~$n$, o complemento
		\[
		C(n,i) \;:=\; S(i)-n
		\]
		é sempre par: $S(i)=3(\text{ímpar})$ é ímpar, $n$ é ímpar,
		logo $C=\text{ímpar}-\text{ímpar}=\text{par}$.
	\end{remark}
	
	% ════════════════════════════════════════════════════════════════════════════
	\section{O Invariante Âncora e o Peso Oracle-Free $\Oms$}
	
	\subsection{A Série Singular como Peso Aritmético}
	
	\begin{definition}[Série Singular de Hardy--Littlewood]
		Para $C\in\NN$ par, define-se
		\[
		\Ss(C) \;:=\!\!\prod_{\substack{p\mid C\\p>2}}\!\frac{p-1}{p-2}.
		\]
	\end{definition}
	
	$\Ss(C)$ é computável em tempo $O(\log C)$ a partir da fatoração de~$C$,
	sem qualquer teste de primalidade de outros números.
	
	\begin{proposition}[$\Ss(C)$ é multiplicativo em $C$, não em $n$]
		\label{prop:mult}
		$\Ss(C)$ é uma função multiplicativa do complemento $C=S(i)-n$, não do
		argumento~$n$. Em particular, $\Ss(C)$ \emph{não} herda a cegueira de
		paridade dos crivos de Selberg, que são multiplicativos em~$n$.
	\end{proposition}
	
	\begin{proof}
		A fórmula de $\Ss$ é um produto de Euler sobre os fatores primos ímpares
		de~$C$, que é determinado pela posição de~$n$ na grade (via $C=S(i)-n$),
		não pelos divisores de~$n$.
		Um primo $p_1$ e um semiprimo $p_2$ que ocupam a mesma célula têm
		complementos $C$ com estruturas de fatoração distintas, gerando valores
		$\Ss(C)$ diferentes; ao passo que o crivo de Selberg, ao avaliar~$n$ por
		seus próprios divisores, não consegue distingui-los.
	\end{proof}
	
	\subsection{Definição de $\Oms$}
	
	\begin{definition}[Peso oracle-free]
		\label{def:omegastar}
		Para cada célula $n=f(r,c)$ em $\Grid$, seja
		$W(n)=\{i: n\in W_i\}$ o conjunto de janelas que contêm~$n$.
		Define-se
		\[
		\Oms(n) \;:=\; \frac{1}{|W(n)|}\sum_{i\in W(n)}\Ss\!\bigl(C(n,i)\bigr),
		\qquad C(n,i)=S(i)-n.
		\]
	\end{definition}
	
	$\Oms(n)$ é computável sem oráculo de primalidade:
	$S(i)$ depende apenas de $N$ e~$i$ (posição);
	$n$ é conhecido por posição na grade;
	$\Ss(C)$ requer apenas a fatoração de $C$ (Observação~\ref{rem:paridade}
	garante que $C$ é par).
	
	\subsection{Conexão com o Estimador de Goldbach}
	
	Define-se o estimador oracle-free para o número de partições de Goldbach
	de um par $a$ como
	\begin{equation}
		\hat\pi_2(a,N) \;:=\!\!\sum_{\substack{p<a/2\\a-p\in\Grid}}
		\Oms(p)\cdot\Oms(a-p).
		\label{eq:estimador}
	\end{equation}
	
	A conjectura de trabalho é que $\hat\pi_2(a,N)\geq c>0$ para $N$
	suficientemente grande em função de~$a$, o que implicaria rigorosamente
	$\pi_2(a)\geq1$.
	A lacuna --- provar essa cota inferior uniformemente em~$a$ ---
	é precisamente a Conjectura de Goldbach em linguagem do crivo geométrico.
	
	% ════════════════════════════════════════════════════════════════════════════
	\section{Verificação Computacional}
	
	\subsection{Metodologia}
	
	Implementamos $\Grid$ computacionalmente para $N$ variando de $10^3$ a
	$10^5$. Para cada célula~$n$ calculamos $\Om(n)$ (oracle-based, contagem
	de vizinhos primos) e $\Oms(n)$ (oracle-free, média de $\Ss(C)$).
	Classificamos $n$ como $P_1$ (primo), $P_2$ (produto de exatamente dois
	fatores primos com multiplicidade) ou composto via crivo de Eratóstenes.
	Aplicamos os testes não-paramétricos KS e MW para comparar as distribuições
	de cada peso sobre $P_1$ e $P_2$.
	O código-fonte está disponível em \cite{bandeira2026code}.
	
	\subsection{Resultados}
	
	\begin{table}[H]
		\centering
		\caption{Comparação dos testes KS ($p<0{,}05$ = significativo) para
			$\Om$ e $\Oms$ em $\Grid$.
			``Sig.'' indica rejeição da hipótese nula de igualdade das distribuições de
			$P_1$ e $P_2$.}
		\label{tab:resultados}
		\smallskip
		\begin{tabular}{%
				r
				l
				r@{\hspace{6pt}}c
				r@{\hspace{6pt}}c
				l
			}
			\toprule
			$N$ & Fatoração &
			\multicolumn{2}{c}{$\Om$ (oracle)} &
			\multicolumn{2}{c}{$\Oms$ (oracle-free)} &
			Resultado \\
			& & $p$-val KS & Sig. & $p$-val KS & Sig. & \\
			\midrule
			$1\,000$   & $2^3\cdot5^3$                    & $1.69\times10^{-3}$  & Sim & $2.57\times10^{-38}$  & Sim & Ambos       \\
			$1\,800$   & $2^3\cdot3^2\cdot5^2$            & $1.18\times10^{-1}$  & Não & $3.15\times10^{-47}$  & Sim & Só $\Oms$   \\
			$2\,400$   & $2^5\cdot3\cdot5^2$              & $1.51\times10^{-1}$  & Não & $3.88\times10^{-57}$  & Sim & Só $\Oms$   \\
			$3\,600$   & $2^4\cdot3^2\cdot5^2$            & $1.30\times10^{-1}$  & Não & $7.70\times10^{-68}$  & Sim & Só $\Oms$   \\
			$7\,200$   & $2^5\cdot3^2\cdot5^2$            & $8.00\times10^{-1}$  & Não & $2.61\times10^{-110}$ & Sim & Só $\Oms$   \\
			$12\,000$  & $2^5\cdot3\cdot5^3$              & $2.54\times10^{-1}$  & Não & $6.75\times10^{-173}$ & Sim & Só $\Oms$   \\
			$15\,000$  & $2^3\cdot3\cdot5^4$              & $1.97\times10^{-5}$  & Sim & $<2\times10^{-224}$   & Sim & Ambos       \\
			$21\,000$  & $2^3\cdot3\cdot5^3\cdot7$        & $6.79\times10^{-23}$ & Sim & $7.02\times10^{-305}$ & Sim & Ambos       \\
			$100\,000$ & $2^5\cdot5^5$                    & $4.04\times10^{-41}$ & Sim & $\ll10^{-100}$        & Sim & Ambos       \\
			\bottomrule
		\end{tabular}
	\end{table}
	
	\subsection{Resultado Central}
	
	\begin{remark}[Superioridade de $\Oms$]
		\label{rem:superior}
		$\Oms$ discrimina $P_1$ de $P_2$ em \textbf{todos} os $N$ testados,
		incluindo $N\in\{1800,2400,3600,7200,12000\}$ cuja fatoração envolve apenas
		$\{2,3,5\}$ e para os quais $\Om$ não detecta diferença significativa.
	\end{remark}
	
	Para $N$ com fatoração restrita a pequenos primos, a grade induz padrões de
	congruência que homogeneízam a distribuição de primos nas vizinhanças ---
	obscurecendo o sinal em $\Om$.
	O complemento $C(n,i)=S(i)-n$, porém, tem fatoração determinada pela
	posição e não pelas obstruções modulares de~$N$.
	Assim $\Ss(C)$ captura um sinal aritmético independente dessas obstruções.
	
	\begin{remark}[Independência dos sinais]
		A correlação entre $\Om$ e $\Oms$ é próxima de zero
		($r\approx-0{,}06$ para $P_1$, $r\approx0{,}02$ para $P_2$),
		confirmando que os dois pesos capturam sinais aritmeticamente independentes.
		$\Oms$ não é uma aproximação de $\Om$: é um objeto distinto.
	\end{remark}
	
	\begin{remark}[Direção do efeito]
		Em todos os $N$ testados, a ordem observada é
		$\Oms(P_1)<\Oms(P_2)<\Oms(\text{composto})$.
		Primos tendem a ter complementos $C$ com mais fatores primos ímpares,
		refletindo a raridade crescente de primos e a estrutura assimétrica do
		complemento aditivo.
		Para a viabilidade do crivo geométrico, o que importa é que as distribuições
		de $\Oms$ sobre $P_1$ e $P_2$ sejam \emph{estatisticamente distinguíveis},
		e é exatamente isso que os testes confirmam.
	\end{remark}
	
	% ════════════════════════════════════════════════════════════════════════════
	\section{Implicações para a Conjectura de Goldbach}
	
	\subsection{A Lacuna Deslocada}
	
	O framework dos três papers pode ser resumido como uma progressão de
	reduções da lacuna central:
	
	\begin{itemize}
		\item \textbf{\cite{bandeira2026unif}:} a lacuna é provar que
		$a-p\in Q_p$ para algum $p$ e todo $a$ específico ---
		equivalente a Goldbach (Teorema~3.5 de \cite{bandeira2026unif}).
		
		\item \textbf{\cite{bandeira2026soma}:} o Lema da Soma Constante
		transforma cada trio primo em testemunha simultânea de três instâncias
		de Goldbach binária e uma de Goldbach ternária, sem reduzir a lacuna,
		mas unificando a linguagem.
		
		\item \textbf{Este paper:} $\Oms$ \emph{desloca} a lacuna.
		Em vez de provar algo sobre a primalidade de~$n$
		(hard: requer distinguir $P_1$ de $P_2$),
		basta provar algo sobre a estrutura aditiva de $C=S(i)-n$.
		Formalmente:
		\[
		\text{Goldbach}(a)
		\;\Leftrightarrow\;
		\exists\,i,N:\;C(p,i)\text{ é par e é soma de dois primos.}
		\]
		Mas $C$ é um número par determinístico, computável por posição.
	\end{itemize}
	
	\subsection{Conexão com o Método do Círculo}
	
	A estimativa de Hardy--Littlewood para o número de representações de
	Goldbach de um par $M$ é $\mathfrak{G}(M)\sim 2C_2\,\Ss(M)\,M/(\ln M)^2$.
	O fator $\Ss(M)$ é exatamente o núcleo de $\Oms$: ao calcular
	$\Ss(C(p,i))$ estimamos, para cada par $(p,a{-}p)$ candidato, quão
	``Goldbach-favorável'' é o complemento~$C$.
	Isso conecta o crivo geométrico diretamente com a heurística de
	Hardy--Littlewood, sem passar pela estrutura multiplicativa do crivo de
	Selberg.
	
	\subsection{O que ainda falta}
	
	A Observação~\ref{rem:superior} mostra que $\Oms$ discrimina $P_1$ de $P_2$
	estatisticamente.
	Isso é necessário mas não suficiente para provar Goldbach.
	A lacuna remanescente é dupla:
	\begin{enumerate}
		\item Provar que $\hat\pi_2(a,N)\geq1$ para todo $a$ específico e $N$
		suficientemente grande, sem usar o oráculo de primalidade para controlar
		os termos de erro.
		\item Garantir que a convergência de $\Oms$ --- a média de $\Ss(C)$ ---
		seja uniforme em~$a$, não apenas em média sobre todos os pares.
	\end{enumerate}
	
	A tensão entre densidade assintoticamente alta (o que os dados confirmam)
	e cobertura pontual de todo~$a$ (o que Goldbach exige) permanece a barreira
	estrutural, agora expressa em termos de $\Ss(C)$ em vez de primalidade de
	vizinhos.
	
	% ════════════════════════════════════════════════════════════════════════════
	\section{Extensões a Outros Problemas Clássicos}
	
	A estrutura do crivo geométrico é agnóstica ao problema-alvo.
	A grade $\Grid$ pode ser ancorada em diferentes regiões da reta numérica,
	e o mesmo aparato --- $S(i)$, $C(n,i)$, $\Ss(C)$, $\Oms$ --- se aplica.
	
	\subsection{Conjectura de Legendre}
	
	A Conjectura de Legendre afirma que existe pelo menos um primo em
	$(n^2,(n{+}1)^2)$ para todo $n\geq1$.
	Ancorada em $L_n$, a grade $G_{3,n}$ tem suas três linhas coincidindo com
	três intervalos de Legendre consecutivos; a linha mediana contém exatamente
	os ímpares de $(n^2,(n{+}1)^2)$.
	O complemento $C(m,i)=S(i)-m$ com $m$ na linha mediana satisfaz
	$C\approx2m\approx2n^2$, e $\Oms$ restrito à linha mediana é um estimador
	oracle-free para a densidade de primos em $(n^2,(n{+}1)^2)$.
	
	\subsection{Primos Gêmeos}
	
	Um par primo gêmeo $(p,p{+}2)$ corresponde a dois elementos adjacentes na
	mesma linha da grade com diferença~2.
	Para uma grade $2\times N$, a configuração diagonal $D_1$ é exatamente um
	par gêmeo.
	O complemento $C(p,i)$ é oracle-free e $\Ss(C)$ estima a densidade de
	gêmeos na vizinhança de~$p$ sem oráculo.
	A conjectura dos primos gêmeos equivale a mostrar que a série de $\Oms$
	sobre as diagonais $D_1$ da grade $2\times N$ diverge --- análogo ao
	argumento de Borel--Cantelli de \cite{bandeira2026unif}.
	
	\subsection{Conjectura de Oppermann}
	
	A Conjectura de Oppermann afirma a existência de primos nos intervalos
	$(n^2,n(n{+}1))$ e $(n(n{+}1),(n{+}1)^2)$ para todo $n$.
	Isso equivale a dois primos na linha mediana do grid de Legendre em posições
	simétricas em relação à âncora $n(n{+}1)=S(i)/3$.
	O complemento $C(m,i)$ para $m$ em cada metade é oracle-free e carrega o
	sinal de Hardy--Littlewood para essa sub-região.
	
	\subsection{Conjectura de Brocard}
	
	A Conjectura de Brocard afirma que existem pelo menos 4 primos entre
	$p_n^2$ e $p_{n+1}^2$.
	A grade ancorada em $p_n^2$ tem linha mediana de comprimento
	$\approx 2p_n\cdot\mathrm{gap}(n)$ ímpares.
	O estimador oracle-free para o número de primos na linha mediana é da ordem
	$p_n\cdot\mathrm{gap}(n)/(2\ln p_n)$, que cresce rapidamente ---
	Brocard pede apenas~4.
	
	\subsection{Atlas de Ancoragens}
	
	\begin{table}[H]
		\centering
		\caption{Instâncias do framework de crivo geométrico para
			diferentes problemas sobre primos.}
		\label{tab:atlas}
		\smallskip
		\begin{tabular}{llll}
			\toprule
			Problema & Ancoragem de $\Grid$ & $K_u$ esperado & Lacuna específica \\
			\midrule
			Goldbach       & Linha inteira                    & $N/(\ln N)^3$                    & Todo par $2M$       \\
			Legendre       & Intervalo $(n^2,\dots)$          & $n/(\ln n)^3$                    & Todo $n$            \\
			Gêmeos         & Grade $2\times N$                & $N/(\ln N)^2$                    & Todo $n$            \\
			Oppermann      & Meia-linha mediana               & $n/(2\ln n)^3$                   & Todo $n$            \\
			Brocard        & Entre $p_n^2$ e $p_{n+1}^2$     & $p_n\cdot\mathrm{gap}/(\ln p)^3$ & Pelo menos 4        \\
			\bottomrule
		\end{tabular}
	\end{table}
	
	% ════════════════════════════════════════════════════════════════════════════
	\section{Investigações Futuras}
	
	\subsection{Formalização Ergódica de $\Oms$}
	
	A Seção~10 de \cite{bandeira2026unif} formula a existência de configurações
	primas como problema de recorrência múltipla no sistema simbólico
	$(X,\mu_N,T)$.
	$\Oms$ oferece um caminho natural para formalizar isso sem hipótese de
	Cramér: o limite de $\Ss(C(n,i))$ quando $N$ cresce pode ser expresso como
	uma média ergódica sobre translações da grade, conectando com os teoremas
	de recorrência de Furstenberg--Katznelson.
	A questão central seria: o limite ergódico de $\Oms$ sobre $P_1$ e $P_2$
	converge para distribuições distintas de forma incondicional, sem hipóteses
	de quasi-independência?
	
	\subsection{$\Oms$ como Peso de Crivo Geométrico}
	
	O programa do crivo geométrico (Seção~11 de \cite{bandeira2026unif}) propõe
	substituir os pesos $\lambda_d$ multiplicativos do crivo de Selberg por
	pesos não-multiplicativos da grade.
	$\Oms$ é o primeiro candidato concreto oracle-free com discriminação
	empiricamente confirmada.
	A investigação natural é: é possível construir um estimador
	$\hat\pi_2(a,N)$ baseado em $\Oms$ para o qual uma cota inferior
	incondicional $\hat\pi_2(a,N)\geq1$ seja provável por métodos de soma de
	exponenciais?
	
	\subsection{Refinamento via Funções $L$ de Dirichlet}
	
	Uma extensão natural de $\Ss(C)$ incorpora caracteres de Dirichlet:
	\[
	\Ss_\chi(C) \;=\;\sum_{\chi\bmod q}\chi(C)\,\frac{L(1,\chi)}{L(2,\chi)}.
	\]
	Esse refinamento capturaria correlações de longa escala entre o complemento
	$C$ e progressões aritméticas, potencialmente controlando os termos de erro
	que impedem a prova de Goldbach pelo método do círculo no caso binário.
	
	\subsection{Extensão a Grades $k\times N$}
	
	Para $k=4$, a soma constante é da forma $4(N+2i)$ e os complementos são
	sempre pares, admitindo $\Ss(C)$.
	A grade $4\times N$ produziria estimadores para Goldbach com
	$K_u\sim N/(\ln N)^4$ --- mais raro que trios, mas potencialmente mais
	tratável analiticamente por ter estrutura simétrica adicional.
	
	\subsection{Testes Computacionais Prioritários}
	
	\begin{enumerate}[(a)]
		\item Rodar $\Oms$ para $N=30030=2\cdot3\cdot5\cdot7\cdot11\cdot13$
		(primorial $p_6\#$): caso onde $\Om$ tinha separação máxima em
		\cite{bandeira2026unif}.
		Esperamos separação de $\Oms$ com estrutura de $\Ss(C)$ distinta,
		revelando qual fator domina.
		
		\item Comparar as distribuições de $C(p,i)$ para $p$ primo vs.\ $p$
		semiprimo: se a fatoração média de $C$ diferir sistematicamente entre
		as classes, isso formalizaria a hipótese de que $\Ss(C)$ capta a
		assinatura aritmética da classe de~$n$.
		
		\item Testar o estimador $\hat\pi_2(a,N)$ (eq.~\eqref{eq:estimador})
		para pares $a$ específicos e comparar com o número real de
		representações de Goldbach: se $\hat\pi_2(a,N)>0$ sempre, isso seria
		evidência computacional forte para a conjectura de trabalho.
	\end{enumerate}
	
	% ════════════════════════════════════════════════════════════════════════════
	\section{Conclusão}
	
	Este trabalho resolve o nó computacional central do programa do crivo
	geométrico iniciado em \cite{bandeira2026unif}: a dependência de $\Om(n)$
	de um oráculo de primalidade.
	A solução usa o Invariante Âncora demonstrado em \cite{bandeira2026soma} ---
	a propriedade $S(i)/3 = f(2,i{+}1)$ --- para construir um complemento
	oracle-free $C(n,i)=S(i)-n$ e um peso $\Oms(n)$ baseado na Série Singular
	de Hardy--Littlewood sobre~$C$.
	
	O resultado computacional é mais forte do que o esperado:
	$\Oms$ não apenas substitui $\Om$ --- ele o supera em robustez,
	discriminando $P_1$ de $P_2$ em todos os $N$ testados, incluindo os $N$
	com obstruções modulares onde o peso original falha.
	A correlação próxima de zero entre os dois pesos confirma que eles capturam
	sinais aritmeticamente independentes.
	
	A lacuna central permanece: provar que $\hat\pi_2(a,N)\geq1$ para todo $a$
	específico exige controlar $\Ss(C)$ uniformemente, reconectando ao problema
	da paridade --- agora em~$C$, não em~$n$.
	Esse deslocamento é o avanço estrutural: o problema é reformulado em termos
	aditivos (estrutura de $C$) em vez de multiplicativos (divisores de $n$),
	abrindo a possibilidade de ataque pelo método do círculo ou por métodos
	ergódicos.
	
	Os três papers formam um framework coerente com Goldbach como instância
	central e outros problemas clássicos --- Legendre, gêmeos, Oppermann,
	Brocard --- como instâncias do mesmo mecanismo geométrico, cada uma com sua
	ancoragem própria e a mesma lacuna estrutural.
	
	% ════════════════════════════════════════════════════════════════════════════
	\begin{thebibliography}{9}
		
		\bibitem{bandeira2026unif}
		T.~Bandeira,
		\textit{Uma Abordagem Unificada para a Conjectura de Goldbach:
			Estrutura Combinatória, Saturação Geométrica e o Elo entre Trios e
			Pares Primos},
		Preprint, maio de 2026.
		
		\bibitem{bandeira2026soma}
		T.~Bandeira,
		\textit{Uma propriedade de soma constante na grade $3\times N$
			e as conjecturas de Goldbach},
		Preprint, maio de 2026.
		
		\bibitem{bandeira2026code}
		T.~Bandeira,
		\textit{\texttt{crivo\_geometrico\_v2.py} --- verificação computacional
			de $\Omega_N$ e $\Omega^*_N$},
		\url{https://github.com/tiagobandeira/conjectura-de-goldbach}, 2026.
		
		\bibitem{HardyLittlewood1923}
		G.~H.~Hardy e J.~E.~Littlewood,
		\textit{Some problems of `Partitio Numerorum'; III: On the expression
			of a number as a sum of primes},
		Acta Mathematica, vol.~44, pp.~1--70, 1923.
		
		\bibitem{Selberg1947}
		A.~Selberg,
		\textit{On an elementary method in the theory of primes},
		Norske Vid.\ Selsk.\ Forh.\ (Trondheim), vol.~19, no.~18,
		pp.~64--67, 1947.
		
		\bibitem{Chen1973}
		J.~R.~Chen,
		\textit{On the representation of a large even integer as the sum of
			a prime and the product of at most two primes},
		Sci.\ Sinica, vol.~16, pp.~157--176, 1973.
		
		\bibitem{Helfgott2013}
		H.~A.~Helfgott,
		\textit{Major arcs for Goldbach's theorem},
		arXiv:1305.2897, 2013.
		
		\bibitem{GreenTao2008}
		B.~Green e T.~Tao,
		\textit{The primes contain arbitrarily long arithmetic progressions},
		Annals of Mathematics, vol.~167, pp.~481--547, 2008.
		
		\bibitem{Cramer1936}
		H.~Cramér,
		\textit{On the order of magnitude of the difference between consecutive
			prime numbers},
		Acta Arithmetica, vol.~2, pp.~23--46, 1936.
		
	\end{thebibliography}
	
\end{document}